\documentclass[11pt]{article}
\usepackage{fontspec,amsmath,amssymb,amsthm,hyperref}
\usepackage[margin=1in]{geometry}
\setlength{\emergencystretch}{3em}
\title{Derive the synchronous one-step splitting bound}
\author{LangevinLean mathematical exposition}
\date{}
\begin{document}\maketitle
\noindent Source module: \texttt{ProximalLangevin/SplittingOneStep.lean}. Compile this proof document with XeLaTeX or LuaLaTeX. The accompanying HTML pairs the mathematical argument with the exact Lean statements and records the scope of the formalization.


In a Hilbert \(L^2\) space, let \(f_1\) be \(m\)-strongly convex with \(L\)-Lipschitz gradient, and let the actual Lipschitz map \(T\) satisfy the nonsmooth energy inequality for \(f\). The reference solves the actual stationary smoothed integral equation for \(f_1+e\), where \(e\le f\le e+\delta\).


\section{Keep every local forcing term explicit}


For \(m,h>0\), drift moment \(G\), noise moment factor \(\sigma\), bounded nonsmooth score \(K\) and value gap \(\delta\), define

\[\mathcal F=(1+(mh)^{-1})[Lh(Gh+\sigma\sqrt h)]^2+2h\delta+2h^2GK+h^2K^2.\]

The four terms are the weighted smooth local defect, the envelope value gap, the drift cross term, and the genuine proximal energy remainder.




\textbf{Correspondence with the displayed declarations.}


\begin{itemize}

\item \texttt{splittingCouplingForcing}: Explicit local forcing in the synchronous splitting proof.
\end{itemize}


\noindent Lean declarations: \texttt{splittingCouplingForcing}.

\section{An exact decomposition makes the energy cancel}


Let \(S=X-h\nabla f_1(X)+B_h\), \(v=(X-h\nabla f_1(X))-(Y_0-h\nabla f_1(Y_0))\), \(R=\int_0^h[\nabla f_1(Y_t)-\nabla f_1(Y_0)]dt\), and \(H=\int_0^h\nabla e(Y_t)dt\). The integral equation gives the exact identity \(S-Y_h=v+R+H\). Expand its square using the rearranged identity; integrate the actual proximal energy inequality for \(T(S)\). The cross-energy lemma cancels the resulting expected \(f\)-difference. Sharp smooth-step contraction and weighted Young inequality bound \(\|v+R\|^2\), while the local-defect lemma bounds \(R\). Discarding \(-\|H\|^2\le0\) yields

\[\|T(S)-Y_h\|_2^2\le(1-mh)\|X-Y_0\|_2^2+\mathcal F.\]

Here \(m\le L\) and \(h\le2/(m+L)\) imply \(mh\le1\). Every reference premise is later supplied by the canonical construction.




Set \(a=\nabla f_1\), \(b=\nabla e\), \(S=X-ha(X)+B_h\), \(v=[X-ha(X)]-[Y_0-ha(Y_0)]\), \(R=\int_0^h[a(Y_t)-a(Y_0)]dt\), and \(H=\int_0^h b(Y_t)dt\). Splitting the actual reference integral equation gives the exact identity \(S-Y_h=v+R+H\). The norm rearrangement is \[\|S-Y_h\|^2=\|v+R\|^2+2\langle S-Y_h,H\rangle-\|H\|^2.\]



Apply the genuine proximal energy estimate to \(T(S)\), then control \(\|v+R\|^2\) using the weighted inequality with \(t=mh\). Strong smooth convexity gives \(\|v\|\le(1-mh)\|X-Y_0\|\), and the step restriction proves \(mh\le1\). The local smooth drift estimate bounds \(\|R\|\), while the preceding cross-energy theorem cancels the original nonsmooth potential difference up to \(2h\delta+2h^2GK\). Discarding \(-\|H\|^2\le0\) leaves exactly the four displayed terms in \texttt{splittingCouplingForcing}. The facts \texttt{hdecomp}, \texttt{hweighted}, \texttt{hRweighted} and \texttt{hcross} correspond to these four stages.




\textbf{Correspondence with the displayed declarations.}


\begin{itemize}

\item \texttt{splitting\_one\_step\_of\_integral\_solution}: One splitting step compared with an actual smooth stationary reference. The pointwise energy premise is proved for the actual averaged joint proximal map; the reference hypotheses are discharged by its canonical Brownian construction. No stochastic error recurrence is assumed.
\end{itemize}


\noindent Lean declarations: \texttt{splitting\_one\_step\_of\_integral\_solution}.

\end{document}
